% TikZ 图：伏安法内外接误差
% Day 28：实验专题
\begin{tikzpicture}[scale=0.86,>=Latex,line cap=round,line join=round]
  \begin{scope}[xshift=-3.25cm]
    \node at (0,1.65) {电流表内接};
    \coordinate (L) at (-1.75,0.75);
    \coordinate (R) at (1.75,0.75);
    \draw[thick] (L) -- (-0.70,0.75);
    \draw[thick] (-0.70,1.03) rectangle (0.05,0.47);
    \node at (-0.325,0.75) {$R_x$};
    \draw[thick] (0.05,0.75) -- (0.70,0.75);
    \draw[thick] (1.00,0.75) circle (0.30);
    \node at (1.00,0.75) {$A$};
    \draw[thick] (1.30,0.75) -- (R) -- (1.75,-0.75) -- (-1.75,-0.75) -- (L);
    % Voltmeter spans both R_x and A in the ammeter-inside connection.
    \draw[thick] (-0.70,0.75) -- (-0.70,-0.05) -- (0.70,-0.05);
    \draw[thick] (1.30,0.75) -- (1.30,-0.05) -- (1.00,-0.05);
    \draw[thick] (0.85,-0.05) circle (0.25);
    \node at (0.85,-0.05) {$V$};
    \node[font=\small,align=center] at (0,-1.35)
      {$V$ 测 $R_x+A$ 的电压\\测得电阻偏大，适合较大电阻};
  \end{scope}
  \begin{scope}[xshift=3.25cm]
    \node at (0,1.65) {电流表外接};
    \draw[thick] (-1.75,0.75) -- (-0.65,0.75);
    \draw[thick] (-0.65,1.03) rectangle (0.15,0.47);
    \node at (-0.25,0.75) {$R_x$};
    \draw[thick] (0.15,0.75) -- (0.70,0.75);
    \draw[thick] (1.00,0.75) circle (0.30);
    \node at (1.00,0.75) {$A$};
    \draw[thick] (1.30,0.75) -- (1.75,0.75) -- (1.75,-0.75) -- (-1.75,-0.75) -- (-1.75,0.75);
    % Voltmeter only spans R_x, so the ammeter reads current through R_x plus the voltmeter branch.
    \draw[thick] (-0.65,0.75) -- (-0.65,-0.05) -- (-0.50,-0.05);
    \draw[thick] (0.15,0.75) -- (0.15,-0.05) -- (0.00,-0.05);
    \draw[thick] (-0.25,-0.05) circle (0.25);
    \node at (-0.25,-0.05) {$V$};
    \node[font=\small,align=center] at (0,-1.35)
      {$A$ 测两支路总电流\\测得电阻偏小，适合较小电阻};
  \end{scope}
\end{tikzpicture}
